\documentclass[10pt,a4paper]{amsart}
\usepackage[margin=19mm]{geometry}
\usepackage{amsmath,amssymb,mathtools}
\usepackage[scr=rsfs,cal=euler]{mathalfa}
\usepackage{xfrac}
\usepackage[unicode,hidelinks,bookmarks=false]{hyperref}
\newcommand{\ie}{i.e.\ }
\newcommand{\cf}{cf.\ }
\newcommand{\resp}{respectively\ }
\newcommand{\Subsection}{\subsection}
\providecommand{\nobreakdash}{\nobreak\hbox{-}}
\setlength{\emergencystretch}{2em}
\multlinegap0pt
\usepackage{bm}
\usepackage{bbm}
\usepackage{dsfont}
\usepackage{xspace}
\usepackage{cancel}
\usepackage{mathtools}
\usepackage{amsthm}
\makeatletter
\@ifundefined{theorem}{\newtheorem{theorem}{Theorem}}{}
\makeatother
\title[Weakly Compatible Mappings and Common Fixed Points Under Gen]{Weakly Compatible Mappings and Common Fixed Points Under Generalized Contractive Conditions}
\author{Alemayehu Negash, Meaza Bogale}
\date{Source: arXiv:2507.00035v1 (CC BY 4.0)}
\begin{document}
\maketitle

\noindent\emph{Source and licence.} Weakly Compatible Mappings and Common Fixed Points Under Generalized Contractive Conditions. Version arXiv:2507.00035v1 (CC BY 4.0), distributed under the Creative Commons Attribution 4.0 International licence, as recorded in the arXiv licence metadata for this exact version and shown on the abstract page https://arxiv.org/abs/2507.00035v1. Full rights notice: licenses/arxiv-2507.00035-CC-BY-4.0.txt.\par\smallskip

\noindent\textit{Editorial scope.} Counted unit: the Theorem whose source label is \texttt{Thrm 3.1} in the article (source0.tex lines 261--279, source label \texttt{Thrm 3.1}), delivered together with the article's own complete original proof of it (source0.tex lines 281--422, one \texttt{proof} environment of 142 lines). No mathematics is written, repaired or re-derived: statement and proof are the article's own text line for line. Required context retained uncounted: none. Every definition, display and label of those lines is retained unchanged. Omitted from this package and not redistributed: 1--260, 280--280, 423--856. Nothing inside a retained excerpt is omitted or altered: every retained body is delivered exactly as the article's lines are written, so no omission receipt of any kind accompanies this package. Citations: the retained excerpts carry no citation command at all, so no bibliography entry of the article is transcribed and no reference is redistributed. Reference adaptation: none was needed and none was made; every reference inside the delivered unit resolves inside it, so no unresolved reference or citation is produced. Printed slips kept literal: no printed word, symbol, hypothesis or number of the counted statement or of its proof was altered. Layout and class: the article's own class is \texttt{article} with packages including amsmath, amssymb, graphicx, amsthm, mathrsfs, mathptmx, mathtools, geometry, enumitem, hyperref; the wrapper is the lane's pinned \texttt{amsart} editorial wrapper at 10pt on A4, which restates the theorem-like environments the retained text needs and adds the article's own macro definitions that the retained text needs (none), with extra wrapper packages (bm, bbm, dsfont, xspace, cancel, mathtools). The delivered numbering is the unit's own. Assets: no retained span contains a figure environment, a graphics-inclusion command, a PDF-page inclusion, a table or a TikZ picture; no image file is copied into this package, the package declares no asset dependency and no image file is redistributed with it. Licence: version arXiv:2507.00035v1 of this article is distributed under the Creative Commons Attribution 4.0 International licence, as recorded in the arXiv licence metadata for this exact version and shown on the abstract page https://arxiv.org/abs/2507.00035v1. A full rights notice is delivered at licenses/arxiv-2507.00035-CC-BY-4.0.txt. Characters: every retained excerpt is pure ASCII, so the delivered engine, XeLaTeX with its default Latin Modern text font, reports no missing character.
\begin{theorem} \label{Thrm 3.1}
Let $K$ be a subset of a metric space $(X,d)$ and $T$, $f$, $g: K \to K$ be self-maps satisfying $\overline{T(K)} \subseteq f(K) \cap g(K)$. Assume:
\begin{multline} \label{(3.1.1)}
d(Tx,Ty) \leq \psi\Bigg(\min\Bigg\{\max\Big\{d(fx,gy),d(Tx,fx),d(Ty,gy),\\
\tfrac{1}{2}\left[d(Tx,gy) + d(Ty,fx)\right]\Big\}, \\
\max\Big\{d(fy,gx),d(Tx,gx),d(Ty,fy),\\
\tfrac{1}{2}\left[d(Tx,fy) + d(Ty,gx)\right]\Big\}\Bigg\}\Bigg)
\end{multline}
for all $x,y \in K$, where $\psi: \mathbb{R}_{+} \to \mathbb{R}_{+}$ is a monotonically increasing continuous function with $0 \leq \psi(t) < t$ for all $t > 0$. If either:
\begin{enumerate}
    \item[(i)] $\overline{T(K)}$ is complete, or
    \item[(ii)] $f(K)$ is complete, or
    \item[(iii)] $g(K)$ is complete
\end{enumerate}
and the pairs $(T,f)$, $(T,g)$ are weakly compatible, then there exists a unique $z \in K$ such that
\[
\{z\} = F(T) \cap F(f) \cap F(g).
\]
\end{theorem}

\begin{proof}
Let $x_0 \in K$. Since $\overline{T(K)} \subset f(K) \cap g(K)$, we can find a sequence $\{x_n\}$ in $K$ such that
\begin{equation}
y_{2n} = Tx_{2n} = fx_{2n+1} \quad \text{and} \quad y_{2n+1} = Tx_{2n+1} = gx_{2n+2}, \quad n=0,1,2,\ldots
\end{equation}

We write
\begin{equation}
\alpha_n = d(y_n, y_{n+1}), \quad n=0,1,2,\ldots
\end{equation}

In the sequence $\{y_n\}$, if $y_{2n} = y_{2n+1}$ for some $n$, then $fx_{2n+1} = Tx_{2n} = Tx_{2n+1} = gx_{2n+2}$. Using this identity and \eqref{(3.1.1)}, it follows that $y_{2n+2} = y_{2n+1}$. Thus $y_{2n+2} = y_{2n+1} = y_{2n}$, and in general, we have $y_m = y_n$ for all $m \geq n$, so that $\{y_m\}_{m \geq n}$ is a constant sequence and hence Cauchy. Therefore, without loss of generality, we assume that $y_n \neq y_{n+1}$ for all $n \in \mathbb{N}$.

Applying \eqref{(3.1.1)} to $x = x_{2n+1}, y = x_{2n}$:
\begin{align*}
\alpha_{2n} &= d(T x_{2n+1}, T x_{2n}) \\
&\leq \psi\Big(\min\Big\{\max\big\{d(fx_{2n+1}, gx_{2n}), d(Tx_{2n+1}, fx_{2n+1}), \\
&\quad d(Tx_{2n}, gx_{2n}), \tfrac{1}{2}[d(Tx_{2n+1}, gx_{2n}) + d(Tx_{2n}, fx_{2n+1})]\big\}, \\
&\quad \max\big\{d(gx_{2n+1}, fx_{2n}), d(Tx_{2n}, fx_{2n}), d(Tx_{2n+1}, gx_{2n+1}), \\
&\quad \tfrac{1}{2}[d(Tx_{2n}, gx_{2n+1}) + d(Tx_{2n+1}, fx_{2n})]\big\}\Big\}\Big) \\
&\leq \psi \Big( \max \big\{ d(f x_{2n+1}, g x_{2n}), d(T x_{2n+1}, f x_{2n+1}), \\
&\quad d(T x_{2n}, g x_{2n}), \tfrac{1}{2}[d(T x_{2n+1}, g x_{2n}) + d(T x_{2n}, f x_{2n+1})] \big\} \Big)
\end{align*}
Substituting $f x_{2n+1} = y_{2n}$, $g x_{2n} = y_{2n-1}$ (for $n \geq 1$), $T x_{2n+1} = y_{2n+1}$, $T x_{2n} = y_{2n}$:
\begin{align*}
d(f x_{2n+1}, g x_{2n}) &= d(y_{2n}, y_{2n-1}) = \alpha_{2n-1} \\
d(T x_{2n+1}, f x_{2n+1}) &= d(y_{2n+1}, y_{2n}) = \alpha_{2n} \\
d(T x_{2n}, g x_{2n}) &= d(y_{2n}, y_{2n-1}) = \alpha_{2n-1} \\
\tfrac{1}{2}[d(T x_{2n+1}, g x_{2n}) + d(T x_{2n}, f x_{2n+1})] &= \tfrac{1}{2}[d(y_{2n+1}, y_{2n-1}) + d(y_{2n}, y_{2n})] \\
&= \tfrac{1}{2} d(y_{2n+1}, y_{2n-1}) \leq \tfrac{1}{2} (\alpha_{2n} + \alpha_{2n-1})
\end{align*}
Thus:
\[
\alpha_{2n} \leq \psi \left( \max \left\{ \alpha_{2n-1}, \alpha_{2n}, \tfrac{1}{2}(\alpha_{2n} + \alpha_{2n-1}) \right\} \right) \leq \psi \left( \max \{\alpha_{2n-1}, \alpha_{2n}\} \right)
\]
If $\max\{\alpha_{2n-1}, \alpha_{2n}\} = \alpha_{2n}$, then $\alpha_{2n} \leq \psi(\alpha_{2n}) < \alpha_{2n}$, contradiction. Hence:
\begin{equation} \label{eq:even-dec}
\alpha_{2n} \leq \psi(\alpha_{2n-1}) < \alpha_{2n-1}
\end{equation}
Similarly, applying \eqref{(3.1.1)} to $x = x_{2n+2}, y = x_{2n+1}$:
\begin{equation} \label{eq:odd-dec}
\alpha_{2n+1} \leq \psi(\alpha_{2n}) < \alpha_{2n}
\end{equation}
Thus $\{\alpha_n\}$ is decreasing and bounded below, so converges to some $r \geq 0$. If $r > 0$, continuity of $\psi$ gives $r \leq \psi(r) < r$, contradiction. Hence:
\begin{equation} \label{eq:dist-conv}
\lim_{n \to \infty} d(y_n, y_{n+1}) = 0
\end{equation}

Next we show that \textbf{$\{y_n\}$ is Cauchy}.
Suppose $\{y_n\}$ is not Cauchy. Then $\exists \varepsilon > 0$ and subsequences $\{m_k\}, \{n_k\}$ with $m_k > n_k > k$ such that:
\begin{align} 
d(y_{2m_k}, y_{2n_k}) &\geq \varepsilon \label{eq:cauchy1} \\
d(y_{2m_k-2}, y_{2n_k}) &< \varepsilon \label{eq:cauchy2}
\end{align}
By triangle inequality:
\[
\varepsilon \leq d(y_{2m_k}, y_{2n_k}) \leq d(y_{2m_k}, y_{2m_k-1}) + d(y_{2m_k-1}, y_{2m_k-2}) + d(y_{2m_k-2}, y_{2n_k})
\]
Using \eqref{eq:dist-conv} and \eqref{eq:cauchy2}:
\begin{equation} \label{eq:lim1}
\lim_{k \to \infty} d(y_{2m_k}, y_{2n_k}) = \varepsilon
\end{equation}
Similarly:
\begin{align}
\lim_{k \to \infty} d(y_{2m_k}, y_{2n_k+1}) &= \varepsilon \label{eq:lim2} \\
\lim_{k \to \infty} d(y_{2m_k-1}, y_{2n_k}) &= \varepsilon \label{eq:lim3} \\
\lim_{k \to \infty} d(y_{2m_k-1}, y_{2n_k+1}) &= \varepsilon \label{eq:lim4}
\end{align}

Now apply \eqref{(3.1.1)} to $x = x_{2n_k+1}, y = x_{2m_k}$:
\begin{align*}
d(y_{2n_k+1}, y_{2m_k}) &\leq \psi \left( \max \left\{ d(y_{2n_k}, y_{2m_k-1}), \alpha_{2n_k}, \alpha_{2m_k-1}, \tfrac{1}{2} d(y_{2n_k+1}, y_{2m_k-1}) \right\} \right)
\end{align*}
Taking $k \to \infty$ and using \eqref{eq:lim1}-\eqref{eq:lim4}:
\[
\varepsilon \leq \psi \left( \max \left\{ \varepsilon, 0, 0, \tfrac{\varepsilon}{2} \right\} \right) = \psi(\varepsilon) < \varepsilon
\]
Contradiction. Hence $\{y_n\}$ is Cauchy in $\overline{T(K)}$. Consequently, $\{Tx_n\}$ is Cauchy, and hence the sequences $\{fx_{2n+1}\}$ and $\{gx_{2n+2}\}$ are Cauchy sequences in $\overline{T(K)} \subset f(K) \cap g(K)$. Since $\{y_n\} \subseteq T(K) \subseteq \overline{T(K)} \subseteq f(K) \cap g(K)$ and one of $\overline{T(K)}$, $f(K)$, or $g(K)$ is complete, $\exists z \in X$ with $\lim y_n = z$.

\subsection*{Case (i): $\overline{T(K)}$ is complete}
Then $Tx_n \to z$ (say) as $n\to\infty$, $z\in\overline{T(K)}$. By the definition of $\{Tx_n\}$, we have
\[
gx_{2n+2}, fx_{2n+1} \to z \quad \text{as} \quad n\to\infty.
\]

Since $\overline{T(K)} \subset f(K) \cap g(K)$, there exist $u,v \in K$ such that $fu = z = gv$.

From inequality \eqref{(3.1.1)}:
\begin{align*}
d(Tx_{2n+1},Tv) &\leq \psi\Big(\min\Big\{\max\big\{d(fx_{2n+1},gv), d(Tx_{2n+1},fx_{2n+1}), d(Tv,gv), \\
&\quad \tfrac{1}{2}[d(Tx_{2n+1},gv) + d(Tv,fx_{2n+1})]\big\}, \max\big\{d(gx_{2n+1},fv), \\
&\quad d(Tv,fv), d(Tx_{2n+1},gx_{2n+1}), \\
&\quad \tfrac{1}{2}[d(Tv,gx_{2n+1}) + d(Tx_{2n+1},fv)]\big\}\Big\}\Big)\\
&\leq \psi\Big(\max\big\{d(fx_{2n+1},gv), d(Tx_{2n+1},fx_{2n+1}), d(Tv,gv), \\
&\quad \tfrac{1}{2}[d(Tx_{2n+1},gv) + d(Tv,fx_{2n+1})]\big\}\Big).
\end{align*}

Letting $n\to\infty$ and using the continuity of $\psi$:
\begin{align*}
d(z,Tv) &\leq \psi\Big(\max\big\{d(z,gv), d(z,z), d(Tv,gv), \tfrac{1}{2}[d(z,gv) + d(Tv,z)]\big\}\Big) \\
&= \psi\Big(\max\big\{0, d(Tv,z), \tfrac{1}{2}[d(z,z) + d(Tv,z)]\big\}\Big) \\
&= \psi\Big(\max\big\{d(z,z), d(Tv,z), \tfrac{1}{2}d(Tv,z)\big\}\Big) \\
&= \psi(d(Tv,z)) < d(Tv,z),
\end{align*}
a contradiction. Hence $z=Tv$, i.e., $Tv=z=gv$.

Similarly, $Tu=z=fu$.

Since $(T,f)$ and $(T,g)$ are weakly compatible and $Tu=fu=z=gv=Tv$, we have $Tfu=fTu$ and $Tgv=gTv$. Hence
\begin{equation}\label{(3.1.13)}
fz = Tz = gz.
\end{equation}

We claim that $z$ is a common fixed point of $T$, $f$, and $g$.

If $z \neq Tz$, then from \eqref{(3.1.1)}:
\begin{align*}
d(z,Tz) &= d(Tu,Tz) \\
&\leq \psi\Big(\min\Big\{\max\big\{d(fu,gz), d(Tu,fu), d(Tz,gz), \\
&\quad \tfrac{1}{2}[d(Tu,gz) + d(Tz,fu)]\big\}, \max\big\{d(gu,fz), d(Tz,fz), \\
&\quad d(Tu,gu), \tfrac{1}{2}[d(Tz,gu) + d(Tu,fz)]\big\}\Big\}\Big) \\
&\leq \psi\Big(\max\big\{d(fu,gz), d(Tu,fu), d(Tv,gz), \\
&\quad \tfrac{1}{2}[d(Tu,gz) + d(Tz,fu)]\big\}\Big) \\
&= \psi\Big(\max\big\{d(z,Tz), d(z,z), d(z,Tz), \\
&\quad \tfrac{1}{2}[d(z,Tz) + d(Tz,z)]\big\}\Big) \\
&= \psi(d(Tz,z)).
\end{align*}
Thus $Tz=z$.
From \eqref{(3.1.13)}:
\[
z \in F(T) \cap F(f) \cap F(g).
\]
The uniqueness follows from \eqref{(3.1.1)}.

\subsection*{Case (ii): $f(K)$ is complete}
Since $\{y_n\}$ is Cauchy, its subsequence $\{fx_{2n+1}\}$ is a Cauchy sequence in $f(K)$. As $f(K)$ is complete, there exists $z \in f(K)$ such that $fx_{2n+1} \to z$. Since the entire sequence $\{y_n\}$ is Cauchy and contains a convergent subsequence, $y_n \to z$. Moreover, since $\{y_n\} \subseteq T(K)$, we have $z \in \overline{T(K)} \subset f(K) \cap g(K)$. Thus there exist $u,v \in K$ such that $fu = z = gv$. The conclusion now follows exactly as in Case (i).

\subsection*{Case (iii): $g(K)$ is complete}
The proof is symmetric to Case (ii) using the subsequence $\{gx_{2n+2}\}$.

Thus in all cases, $F(T) \cap F(f) \cap F(g) = \{z\}$.
\end{proof}

\end{document}
